% % % % % % % % % % % % % % % % % % % % % % % % % % % % % % % % % % % % % % % % % % % % % % % % % % % % % % % % % % % % % % % %
% SAT_STYLE_MATH_MCQS.tex
% 1_15.tex
% 
% encoding: UTF-8 
% EOL: LF
%
% format: LaTeX
% indent: spaces (2)
% width: 127
% % % % % % % % % % % % % % % % % % % % % % % % % % % % % % % % % % % % % % % % % % % % % % % % % % % % % % % % % % % % % % % %
Soit $\theta$ l'angle géométrique formé par l'aiguille des heures et celle des
minutes. Quel est l'abaissement de $\theta$ entre $12$\,h\,$12$ et
$15$\,h\,$15$?
\begin{enumerate}
  \item $31,5°$
  \item $-31,5°$
  \item $64,5°$
  \item $-58,5°$
  \item $58,5°$
\end{enumerate}
%
%
\begin{proof}
À $12$\,h\,$12$: en $12$ minutes, l'aiguille des heures avance de $6°$ et
celle des minutes de $72°$, donc $\theta_{12:12} = 72° - 6° = 66°$.

À $15$\,h\,$15$: angle initial $90°$. En $15$ minutes, l'aiguille des heures
avance de $7,5°$ et celle des minutes de $90°$, donc
$\theta_{15:15} = 90° + 7,5° - 90° = 7,5°$.

L'abaissement est $\theta_{15:15} - \theta_{12:12} = 7,5° - 66° = -58,5°$. Réponse \textbf{E}.
\end{proof}